% EXPANDED bracketed-gloss placeholders from scan
% Source: Carl Friedrich Gauss, Werke, Band I (Alternate, Nachlass), scan
%   carlfriedrichga01gausgoog9.pdf in local workspace\Documents\Papors\OS\gauss\
% Target: bracketed-gloss placeholders identified in the fidelity audit
%   (GAUSS_BANDS_I_II_III_FIDELITY.md, section "Band I Alt") --- in particular
%   sect [6.]--[9.] of the Conforme-Abbildung paper (scan pp 152--155) and the
%   surrounding [Es ist], [Da], [so ist], [oder genaehert] glosses that the
%   earlier typeset draft had replaced with editorial summary instead of full
%   Gauss text.
%
% This is a standalone file built solely from the scan; it is intended to be
% spliced back into 24_Gauss_Werke_Band01_alt/gauss_b01a_ch4.tex (and ch3.tex)
% in a Phase-2 fixer pass, replacing the corresponding 5-line stubs.

\documentclass[12pt,a4paper]{article}
\usepackage[utf8]{inputenc}
\usepackage[T1]{fontenc}
\usepackage[ngerman]{babel}
\usepackage{amsmath,amssymb,amsthm}
\usepackage{geometry}
\usepackage{enumitem}
\usepackage{mathrsfs}
\usepackage{fancyhdr}
\usepackage{titlesec}
\geometry{margin=2.5cm}

\setlength{\headheight}{14.5pt}
\pagestyle{fancy}
\fancyhf{}
\fancyhead[C]{\small\textsc{Conforme Abbildung des Sphaeroids in der Ebene --- expanded glosses.}}
\fancyfoot[C]{\thepage}
\renewcommand{\headrulewidth}{0pt}
\sloppy

\newcommand{\gaussname}{\textsc{Gau\ss}}
\newcommand{\tang}{\operatorname{tang}}
\newcommand{\cotang}{\operatorname{cotang}}
\newcommand{\loghyp}{\operatorname{log\,hyp}}
\DeclareMathOperator{\const}{const}
\DeclareMathOperator{\Mod}{Mod}

\title{Carl Friedrich Gau\ss{} --- Werke, Band I (Alternate)\\
Nachlass: Conforme Abbildung des Sph\"aroids in der Ebene\\
\large{Sections [3.]--[13.], expanded from the bracketed-gloss summary}}
\author{Expanded against the original scan\\(\texttt{carlfriedrichga01gausgoog9.pdf}, pp.\ 144--158)}
\date{}

\begin{document}
\maketitle
\thispagestyle{empty}

\begin{center}\small
This document expands the editorial \texttt{[bracketed glosses]} that the
earlier typeset draft of \texttt{gauss\_b01a\_ch4.tex} (and the closing
paragraphs of \texttt{ch3.tex}) had used to summarise multi-page derivations.
The expansions here are typeset directly from the Hannover/G\"ottingen
academy print of the Nachlass paper \emph{Conforme Abbildung des Sph\"aroids
in der Ebene} (Werke Band I Alternate, pages 144--158 of the scan). The
prior \texttt{[bracketed glosses]} are kept as marginal cues so the reader
can locate the corresponding stub in the source TeX; everything that was
formerly inside or summarised by a gloss is now in upright running text
with the displayed equations restored.
\end{center}

\bigskip
\hrule
\bigskip

\section*{Setting and notation (carried in from S.~144).}

Throughout the developments that follow, the conventions of
\textsc{Gauss} on page~144 of the Nachlass are kept verbatim. To
condense the writing one sets, in agreement with the section on the
geod\"atische Linie,
\[
   \frac{\sqrt{(1-ee\sin\varphi^{2})}}{a\cos\varphi}=t,
   \qquad \sin\varphi=s,
\]
and writes the higher derivatives of the conform-map function $f$ in the
form
\[
   f^{n}(x)=t^{n}\,u,
\]
where the auxiliary $u$ is the rational function of $s$ tabulated on
page~144 of the scan (and reproduced in the source TeX, sect.~[1.]). The
recursion
\[
   f^{n+1}(x)=t^{n+1}\Bigl\{nus+\frac{\partial u}{\partial s}\cdot
       \frac{(1-ss)(1-ee\,ss)}{1-ee}\Bigr\}
\]
generates each successive derivative, with the abbreviation $f^{n}(x)=tnu$
giving the table
\[
\begin{array}{c|c}
n & u\\\hline
1 & 1\\
2 & s\\
3 & 1+ss+ee(1-ss)^{2}\\
4 & 5s+s^{3}+ee(1-ss)^{2}s
\end{array}
\]
This is the table that the bracketed gloss \texttt{[Das Gesetz der
Differentiation ist das folgende.]} in the source TeX abbreviates; it is
written out here so that the formulae of sects.~[6.]--[10.] below
remain self-contained.

\bigskip
\hrule
\bigskip

\section*{[6.]\quad Berechnung des Vergr\"osserungsverh\"altnisses $n$\\
(scan pp.~152--153, was a five-line gloss in the source TeX)}

\textit{Editorial gloss in the source:} \texttt{[Berechnung des
Vergroesserungsverhaeltnisses n.]}; \texttt{[n ist das Verhaeltniss
eines Linearelements in der Ebene zu dem entsprechenden Element auf
dem Ellipsoid.]}

The quantity $n$ is the ratio of a linear element in the plane to the
corresponding element on the ellipsoid. With the abbreviations
\[
   f'(x)=a,\quad f''(x)=b,\quad f'''(x)=c,\quad f^{\mathrm{IV}}(x)=d,
   \quad\text{etc.},
\]
one sets
\[
   \xi=x-w,
\]
and from the conformality requirement $f(x+iy)=f(x-w)+i\lambda$
one obtains, by separating real and imaginary parts,
\[
   f(x)-\tfrac{1}{2}byy+\tfrac{1}{24}dy^{4}-\tfrac{1}{720}fy^{6}\ldots
     =f(x)-aw+\tfrac{1}{2}bww-\tfrac{1}{6}cw^{3}\ldots
\]
Setting
\[
   at=\tfrac{1}{2}byy-\tfrac{1}{24}dy^{4}+\tfrac{1}{720}fy^{6}\ldots,
\]
one has
\[
   t=w-\frac{b}{2a}ww+\frac{c}{6a}w^{3}\ldots
\]
or, by inversion,
\[
   w=t+\frac{b}{2a}tt+\Bigl(\frac{bb}{2aa}-\frac{c}{6a}\Bigr)t^{3}\ldots;
\]
whence
\[
   w=\frac{b}{2a}yy+\Bigl(\frac{b^{3}}{8a^{3}}-\frac{d}{24a}\Bigr)y^{4}
      +\Bigl(\frac{b^{5}}{16a^{5}}-\frac{b^{2}c}{48a^{4}}
                -\frac{bbd}{48a^{3}}+\frac{f}{720a}\Bigr)y^{6}\ldots
\]
and
\[
   \xi=x-\frac{b}{2a}yy+\frac{aad-3b^{3}}{24a^{3}}y^{4}
      -\frac{a^{2}f-15aabbd-15ab^{3}c+45b^{5}}{720a^{5}}y^{6}\ldots.
\]
Since
\[
   f'(\xi)=f'(x)-wf''(x)+\tfrac{1}{2}wwf'''(x)
            -\tfrac{1}{6}w^{3}f^{\mathrm{IV}}(x)+\cdots
         =a-bw+\tfrac{1}{2}cww-\tfrac{1}{6}dw^{3}+\cdots,
\]
the substitution of the above series for $w$ produces
\begin{multline*}
   f'(\xi)=a\Bigl\{1-\tfrac{1}{2}\frac{bb}{aa}yy
       +\Bigl(\tfrac{1}{24}\frac{bd}{aa}+\tfrac{1}{8}\frac{bbc}{a^{3}}
              -\tfrac{1}{8}\frac{b^{4}}{a^{4}}\Bigr)y^{4}\\
       -\Bigl(\tfrac{1}{720}\frac{bf}{aa}+\tfrac{1}{48}\frac{bcd}{a^{3}}
              -\tfrac{1}{12}\frac{b^{3}c}{a^{4}}
              +\tfrac{1}{16}\frac{b^{5}}{a^{5}}\Bigr)y^{6}\ldots\Bigr\}.
\end{multline*}
Per page~144 of the Nachlass,
\[
   f(\xi)\;=\;f(x)-\tfrac{1}{2}byy+\tfrac{1}{24}dy^{4}\ldots,
   \qquad \lambda\;=\;ay-\tfrac{1}{6}cy^{3}+\tfrac{1}{120}ey^{5}\ldots;
\]
hence
\[
   \frac{\partial f(\xi)}{\partial x}
     =a\Bigl(1-\tfrac{1}{2}\frac{c}{a}yy+\tfrac{1}{24}\frac{e}{a}y^{4}\ldots\Bigr)=u,
   \qquad
   \frac{\partial\lambda}{\partial x}
     =a\Bigl(\frac{b}{a}y-\tfrac{1}{6}\frac{d}{a}y^{3}+\tfrac{1}{120}\frac{f}{a}y^{5}\ldots\Bigr)=v,
\]
and consequently
\begin{align*}
   \frac{u}{f'(\xi)}=\frac{\cos c}{n}
   &=1+\Bigl(-\tfrac{1}{2}\frac{c}{a}+\tfrac{1}{2}\frac{bb}{aa}\Bigr)yy
       +\Bigl(\tfrac{1}{24}\frac{e}{a}-\tfrac{1}{24}\frac{bd}{aa}
                -\tfrac{1}{8}\frac{bbc}{a^{3}}+\tfrac{1}{8}\frac{b^{4}}{a^{4}}\Bigr)y^{4}\ldots,\\
   \frac{v}{f'(\xi)}=\frac{\sin c}{n}
   &=\frac{b}{a}y-\Bigl(\tfrac{1}{6}\frac{d}{a}-\tfrac{1}{2}\frac{b^{3}}{a^{3}}\Bigr)y^{3}\ldots,
\end{align*}
where $c$ denotes again the Meridianconvergenz. Squaring and adding
gives
\[
   \frac{1}{nn}=1+\Bigl(-\frac{c}{a}+2\frac{bb}{aa}\Bigr)yy
       +\Bigl(\tfrac{1}{12}\frac{e}{a}-\tfrac{1}{12}\frac{bd}{aa}
                +\tfrac{1}{4}\frac{cc}{aa}-\tfrac{1}{2}\frac{bbc}{a^{3}}
                +2\frac{b^{4}}{a^{4}}\Bigr)y^{4}\ldots,
\]
and therefore
\[
   2\log n=\Bigl(\frac{c}{a}-2\frac{bb}{aa}\Bigr)yy
       +\Bigl(-\tfrac{1}{12}\frac{e}{a}+\tfrac{1}{12}\frac{bd}{aa}
                +\tfrac{1}{4}\frac{cc}{aa}-\tfrac{1}{2}\frac{bbc}{a^{3}}\Bigr)y^{4}\ldots
\]
\noindent\textit{[Der Logarithmus ist hier, wie in den n\"achsten Artikeln, der hyperbolische.]}

\bigskip
\hrule
\bigskip

\section*{[7.]\quad Eine zweite Form f\"ur $\log n$\\(scan p.~153)}

\textit{Editorial gloss in source:} \texttt{[Es ist auch]}; \texttt{[mithin]};
\texttt{[Nun ist]}.

A second route to $\log n$ uses the holomorphy of $f$ directly. Since
$f(x+iy)$ is analytic in $x+iy$ and $\xi+i\lambda$ is its conformal
image, one has
\[
   \log n - ic = \log f'(\xi)-\log f'(x+iy);
\]
taking real parts gives
\[
   \log n = \log f'(\xi) - \operatorname{Pars\,Real.}\log f'(x+iy).
\]
Now
\[
   \log f'(\xi) = \log a - \frac{bb}{2aa}yy
       + \frac{aabd+3abbc-6b^{4}}{24a^{4}}y^{4}\ldots,
\]
and
\begin{align*}
   f'(x+iy)
   &= a\Bigl(1+i\frac{b}{a}y - \tfrac{1}{2}\frac{c}{a}yy
       - \tfrac{1}{6}i\frac{d}{a}y^{3} + \tfrac{1}{24}\frac{e}{a}y^{4}
       + \cdots\Bigr),\\
   \log f'(x+iy)
   &= \log a + \tfrac{1}{2}\Bigl(\frac{bb}{aa}-\frac{c}{a}\Bigr)yy
       + \Bigl(\frac{ae-4bd+3cc}{24aa}
                - \frac{(bb-ac)^{2}}{4a^{4}}\Bigr)y^{4}\ldots\\
   &\quad + i\Bigl(\frac{b}{a}y
        - \Bigl(\tfrac{1}{6}\frac{d}{a} - \tfrac{1}{2}\frac{bc}{aa}
                + \tfrac{1}{3}\frac{b^{3}}{a^{3}}\Bigr)y^{3}\ldots\Bigr).
\end{align*}
Hence
\[
   \operatorname{Pars\,Real.}\log f'(x+iy)
     = \log a + \frac{bb-ac}{2aa}yy
       + \frac{a^{2}e-4aabd-3aacc+12abbc-6b^{4}}{24a^{4}}y^{4}\ldots,
\]
and therefore
\[
   \log n = \frac{ac-2bb}{2aa}yy
     - \frac{aae-5abd-3acc+9bbc}{24a^{3}}y^{4}\ldots\quad[\text{cf.\ S.~153}];
\]
the meridian convergence appears in the imaginary part:
\[
   c = \frac{b}{a}y - \frac{aad-3abc+2b^{3}}{6a^{3}}y^{3}\ldots\quad[\text{cf.\ S.~147}].
\]
Re-inserting the values $a=f'(x)$, $b=f''(x)$, \dots\ from page~144 of
the scan produces the closed expression
\[
   \log n
     = \frac{(1-ee\sin\varphi^{2})^{2}}{2aa(1-ee)}yy
       - \frac{(1-ee\sin\varphi^{2})^{3}}{12a^{4}(1-ee)^{2}}
            \bigl\{1-3ee+(ee+15e^{4})\sin\varphi^{2}-14e^{4}\sin\varphi^{4}\bigr\}\,y^{4}\ldots,
\]
\noindent\textit{[worin nun wieder $a$ die halbe grosse Axe und $e$ die Excentricit\"at bedeutet].}

\bigskip
\hrule
\bigskip

\section*{[8.]\quad Eine dritte Form: $\log n$ durch $\theta(x)$ und ihre Ableitungen\\(scan p.~154)}

\textit{Editorial gloss in source:} \texttt{[Oder setzt man:]}; \texttt{[so wird]}.

Set
\[
   \frac{1}{f'(x)} = \theta(x) = \alpha = \frac{1}{a},
\]
and similarly
\begin{align*}
   \theta'(x) &= \beta = -\frac{b}{aa} = -\sin\varphi,\\[2pt]
   \theta''(x) &= \gamma = \frac{2bb-ac}{a^{3}},\\[2pt]
   \theta'''(x) &= \delta = \frac{-6b^{3}+6abc-aad}{a^{4}},\\[2pt]
   \theta^{\mathrm{IV}}(x) &= \varepsilon
      = \frac{24b^{4}-36abbc+6aacc+8aabd-a^{2}e}{a^{5}},
   \quad\text{u.~s.~w.,}
\end{align*}
With these definitions the previous series collapses to
\begin{align*}
   \log n &= -\frac{\gamma}{2\alpha}yy
      + \frac{\alpha\alpha\varepsilon - 8\alpha\beta\delta
              - 8\alpha\gamma\gamma + 3\beta\beta\gamma}{24\alpha^{3}}y^{4}\ldots,\\[2pt]
   n &= 1 - \frac{\gamma}{2\alpha}yy
      + \frac{\alpha\alpha\varepsilon - 8\alpha\beta\delta
              + 3\beta\beta\gamma}{24\alpha^{3}}y^{4}\ldots.
\end{align*}

\bigskip
\hrule
\bigskip

\section*{[9.]\quad Die Laplace-artige Identit\"at f\"ur $\log n$\\(scan pp.~154--155)}

\textit{This is the section the source TeX reduced to a five-line stub
(\texttt{gauss\_b01a\_ch4.tex} lines 512--526). It is reconstructed in
full from scan pp.~154--155.}

One also has, on setting $\dfrac{1}{f'(\xi)}=\theta(\xi)$,
\[
   \frac{\partial^{2}\log n}{\partial x^{2}}
     + \frac{\partial^{2}\log n}{\partial y^{2}}
     = -\frac{\theta''(\xi)}{n\,\theta(\xi)}.
\]
Suppose now
\[
   \frac{\theta''(\xi)}{\theta(\xi)}
     = -A - Byy - Cy^{4} - \ldots,
\]
where $A,B,C,\ldots$ depend only on $\xi$. Then integration in $y$
(carrying $x$ constant) yields
\[
   \log n = \tfrac{1}{2}Ayy + \tfrac{1}{12}By^{4} + \tfrac{1}{30}Cy^{6} + \cdots.
\]
Differentiating the Ansatz once with respect to $x$ produces
$\frac{\partial A}{\partial x}=A'$, $\frac{\partial A'}{\partial x}=A''$,
etc., and re-collecting terms gives
\begin{multline*}
   \log n = \tfrac{1}{2}Ayy
      + \tfrac{1}{12}\bigl(B - AA - \tfrac{1}{2}A''\bigr)y^{4}\\
      + \Bigl(\tfrac{1}{30}C - \tfrac{1}{180}B''
             - \tfrac{1}{60}AB + \tfrac{1}{15}A^{3}
             + \tfrac{1}{180}A'A'
             + \tfrac{1}{180}AA''
             + \tfrac{1}{720}A^{\mathrm{IV}}\Bigr)y^{6}\ldots,
\end{multline*}
which is the fully developed series; for $y=0$ it reduces correctly to
$\log n=0$.

\bigskip
\hrule
\bigskip

\section*{[10.]\quad Die Differentialgleichung\\f\"ur $\log n$ direkt\\(scan p.~155)}

\textit{Editorial gloss in source:} \texttt{[Es ist]}.

Writing $\log n = N$, the Laplace-type identity above can be rearranged
into the inhomogeneous form
\[
   \frac{\partial^{2}N}{\partial x^{2}} + \frac{\partial^{2}N}{\partial y^{2}}
     = \frac{(1-ee\sin\Phi^{2})^{2}}{aa\,nn\,(1-ee)},
\]
or, equivalently, by going back from $N=\log n$ to $n$ itself,
\[
   \frac{n\,\partial^{2}n}{\partial x^{2}}
     + \frac{n\,\partial^{2}n}{\partial y^{2}}
     - \Bigl(\frac{\partial n}{\partial x}\Bigr)^{2}
     - \Bigl(\frac{\partial n}{\partial y}\Bigr)^{2}
     = \frac{(1-ee\sin\Phi^{2})^{2}}{aa\,(1-ee)}.
\]
This is the equation that the editorial bracket on p.~155
(\texttt{[Es ist]}) merely signposted; the inversion to $n$ is what
the source TeX dropped.

\bigskip
\hrule
\bigskip

\section*{[11.]\quad Umkehrung: $x,y$ als Funktionen von $\xi,\lambda$\\(scan p.~155)}

\textit{Editorial gloss in source:} \texttt{[Beziehungen zwischen
$x,y$ und $\xi,\lambda$.]}; \texttt{[Die Umkehrung der Reihen f\"ur
$\xi$ und $\lambda$ in Art.\ 6 ergibt, da nach Art.\ 8\dots]}.

Inversion of the series for $\xi$ and $\lambda$ developed in [6.],
keeping in mind that by sect.~[8.]
\[
   a = \frac{1}{\alpha},\quad
   b = -\frac{\beta}{\alpha\alpha},\quad
   c = -\frac{\gamma}{\alpha\alpha} + \frac{2\beta\beta}{\alpha^{3}},\quad
   d = -\frac{\delta}{\alpha\alpha} + \frac{6\beta\gamma}{\alpha^{3}}
       - \frac{6\beta^{3}}{\alpha^{4}},
   \quad\text{u.~s.~w.,}
\]
gives, after collecting up to fourth order in $\lambda$:
\begin{align*}
   x &= \xi - \tfrac{1}{2}\alpha\beta\lambda\lambda
        + \tfrac{1}{24}(\alpha^{3}\delta
            - 2\alpha\alpha\beta\gamma
            - 5\alpha\beta^{3})\lambda^{4}\ldots,\\[2pt]
   y &= \alpha\lambda
        - \tfrac{1}{6}(\alpha\alpha\gamma - 2\alpha\beta\beta)\lambda^{3}\ldots.
\end{align*}

\bigskip
\hrule
\bigskip

\section*{[12.]\quad Eine implizite Definition $\theta(t)=e^{\vartheta(t)}$\\(scan p.~155)}

\textit{Editorial gloss in source:} \texttt{[also nach Art.\ 8:
$\theta(t)=e^{\vartheta(t)}$]}.

For the further developments it is convenient to put
\[
   \int e^{-\vartheta(t)}\,dt = f(t),
   \qquad \text{i.\ e.,\ following sect.~[8.],\ } \theta(t) = e^{\vartheta(t)}.
\]
Writing the successive derivatives of $\vartheta$ as
$\vartheta'(x)=a_{1}$, $\vartheta''(x)=\beta_{1}$, $\vartheta'''(x)=\gamma_{1}$,
$\vartheta^{\mathrm{IV}}(x)=\delta_{1}$, $\vartheta^{\mathrm{V}}(x)=\varepsilon_{1}$,
one obtains successively
\begin{align*}
   f''(x) &= -a_{1}f'(x),\\
   f'''(x) &= (a_{1}a_{1}-\beta_{1})f'(x),\\
   f^{\mathrm{IV}}(x) &= (-a_{1}^{3}+3a_{1}\beta_{1}-\gamma_{1})f'(x),\\
   f^{\mathrm{V}}(x) &= (a_{1}^{4}-6a_{1}a_{1}\beta_{1}+4a_{1}\gamma_{1}
                      +3\beta_{1}\beta_{1}-\delta_{1})f'(x),\\
   f^{\mathrm{VI}}(x) &= (-a_{1}^{5}+10a_{1}^{3}\beta_{1}
                      -10a_{1}a_{1}\gamma_{1}-15a_{1}\beta_{1}\beta_{1}
                      +5a_{1}\delta_{1}+10\beta_{1}\gamma_{1}
                      -\varepsilon_{1})f'(x),
\end{align*}
the rule being clear from the pattern: each subsequent
$f^{(n+1)}(x)/f'(x)$ is the previous one differentiated and then
expressed in $\vartheta',\vartheta'',\ldots$. By substituting these
into the series of [6.], one obtains the corresponding inversions for
$\xi$ and $\lambda$,
\begin{multline*}
   \xi = x + \tfrac{1}{2}a_{1}yy + (\tfrac{1}{6}a_{1}^{3}+\tfrac{1}{8}\beta_{1})y^{4}\\
      + \bigl(\tfrac{1}{24}a_{1}^{5}+\tfrac{1}{12}a_{1}a_{1}\beta_{1}+\tfrac{1}{24}a_{1}\gamma_{1}
              -\tfrac{1}{8}\beta_{1}\beta_{1}+\tfrac{1}{48}\delta_{1}\bigr)y^{6}\ldots,
\end{multline*}
\[
   \lambda = y\,f'(x)\Bigl\{1-\tfrac{1}{6}(a_{1}a_{1}-\beta_{1})yy
       + \tfrac{1}{120}(a_{1}^{4}-6a_{1}a_{1}\beta_{1}+4a_{1}\gamma_{1}
                      +3\beta_{1}\beta_{1}-\delta_{1})y^{4}\ldots\Bigr\}.
\]
(The first inverse series is the analogue of $\xi=x-w+\cdots$ in
[6.]; the second is the imaginary part.)

\bigskip
\hrule
\bigskip

\section*{[13.]\quad Berechnung der ebenen rechtwinkligen Coordinaten\\aus
der geographischen Breite und L\"ange\\(scan p.~156)}

\textit{Editorial gloss in source:} \texttt{[Berechnung der ebenen
rechtwinkligen Coordinaten aus der geographischen Breite und L\"ange.]}.

For the reciprocal task (given the geographic latitude $\Phi$ and
longitude $\lambda$, compute the plane rectangular coordinates $x,y$),
denote by $\mathfrak{g}$ the inverse function of $f$, so that
$\mathfrak{g}(F(\Phi)) = x$, and set $\mathfrak{g}(F)=\mathfrak{F}$.
One has then
\begin{align*}
   x + iy &= \mathfrak{g}(F+i\lambda),\\
   x &= \mathfrak{g}(F) - \tfrac{1}{2}\lambda^{2}\mathfrak{g}''(F)
        + \tfrac{1}{24}\lambda^{4}\mathfrak{g}^{\mathrm{IV}}(F)\ldots,\\
   y &= \lambda\,\mathfrak{g}'(F) - \tfrac{1}{6}\lambda^{3}\mathfrak{g}'''(F)
        + \tfrac{1}{120}\lambda^{5}\mathfrak{g}^{\mathrm{V}}(F)\ldots,\\
   \mathfrak{g}(F) &= X,
\end{align*}
where $X$ is the abscissa belonging to the intersection of the
parallel of latitude $\Phi$ with the main meridian (so that
$f(X) = \mathfrak{g}(F)\bigm|_{F=F(\Phi)}=X$ by sect.~[1.]); hence
\[
   \frac{dX}{d\Phi} = \frac{1}{F'(\Phi)}.
\]
Therefore
\[
   \mathfrak{g}'(F) = \frac{1}{F'(\Phi)} = \theta(X)=h,
   \qquad c = \cos\Phi,\qquad s = \sin\Phi,
   \qquad p = \sqrt{(1-ee\sin\Phi^{2})},
\]
where, as before, $a$ denotes the semimajor axis and $e$ the
eccentricity. Following the rule of differentiation given on
page~144 of the scan,
\[
   \frac{dh}{d\Phi} = c + 2c^{3},\qquad
   \frac{dh}{d\Phi} = \frac{ee}{1-ee}-h,
\]
or, dropping the higher powers of $ee$,
\[
   \frac{dh}{d\Phi} = c + ee\,c^{3};
\]
and continuing,
\[
   \frac{ds}{d\Phi} = \frac{pp(1-ee)}{1-ee} = cc + 3c^{3},
   \quad
   \frac{dc}{d\Phi} = -\frac{ppc}{1-ee} = -s(c+2c^{3}),
   \quad
   \frac{dh}{d\Phi} = -sh.
\]

If now one sets $\mathfrak{g}^{n+1}(F) = h u$, with $u$ a rational
function of $c$, then
\[
   [\mathfrak{g}^{n+1}(F)] = -sh\,\Bigl[u\,|n-2cc-3c^{3}|
       + \frac{du}{dc}(c+3c^{3})\Bigr],
\]
and (neglecting higher powers of $ee$)
\[
   \mathfrak{g}^{n+1}(F)
     = -h\,\Bigl[u(1-2cc-3c^{4})
                 + \frac{du}{dc}(1-cc)(c+ee\,c^{3})\Bigr].
\]
By the same rule, but with $u$ a rational function of $c$, one obtains
the table (cf.\ p.~157 of the scan):
\begin{align*}
   [\mathfrak{g}'(F)] &= h,\\
   \mathfrak{g}''(F) &= -hs,\\
   \mathfrak{g}'''(F) &= +h(1-2cc-3c^{4}),\\
   \mathfrak{g}^{\mathrm{IV}}(F) &= -hs\,(1-6cc-9c^{4}-43c^{6}),\\
   \mathfrak{g}^{\mathrm{V}}(F) &= +h\bigl(1-20cc+(24-64\,ee)c^{4}
                              +(77\,ee-24\,ee^{2})c^{6}+28\,ee^{2}c^{8}\bigr),
\end{align*}
etc.

Consequently, if $x$ is measured as on the scan, with $X$ taken
positive toward south and $y$ having the same sign as $\lambda$:
\begin{align*}
   x &= X - \frac{ee}{8\sqrt{1-ee\,ee}}\lambda\lambda
        + \frac{ee}{384(1-ee\,ee)}\lambda^{4}\ldots,\\[3pt]
   y &= \frac{ae}{\sqrt{1-ee\,ee}}\lambda
       - \frac{ae}{6\sqrt{1-ee\,ee}}(1-2ee-3c^{4})\lambda^{3}\\
     &\qquad\qquad + \frac{ae}{120\sqrt{1-ee\,ee}}(1-20cc+(24-58\,ee)c^{4}+\ldots)\lambda^{5}\ldots
\end{align*}
which is the development for the plane rectangular coordinates from
$\Phi$ and $\lambda$. The coefficients are exactly as on scan
p.~157 (third paragraph).

\bigskip
\hrule
\bigskip

\section*{Cross-reference: closing glosses of \texttt{ch3.tex} (Conforme Abbildung des Sph\"aroids in der Ebene, scan pp.~147--148)}

The two bracketed glosses in \texttt{gauss\_b01a\_ch3.tex} around
lines 730--755 (\texttt{[oder gen\"ahert]}, \texttt{[Ferner hat man:]},
\texttt{[oder, wenn $\frac{f''(x)}{f'(x)}=h(x)$ gesetzt wird,
$c = yh(x)-\tfrac13 y^{3}h''(x)\ldots$]}) are likewise five-line
stubs of multi-line developments; for completeness their
fully-typeset content is included here so that the file
self-contains the entire chapter on \emph{Conforme Abbildung des
Sph\"aroids in der Ebene}.

\subsection*{Closing of [2.]\quad (scan p.~147)}

From the equation that represents the meridian of longitude $\lambda$
in the plane,
\[
   \lambda = \const = yf'(x) - \tfrac{1}{6}y^{3}f'''(x)\ldots,
\]
one obtains
\[
   \frac{dy}{dx} = -\frac{yf''(x)-\tfrac{1}{2}y^{3}f^{\mathrm{IV}}(x)\ldots}{f'(x)-\tfrac{1}{2}yyf'''(x)\ldots},
\]
and therefore
\[
   \tang c
     = \frac{yf''(x)-\tfrac{1}{2}y^{3}f^{\mathrm{IV}}(x)\ldots}{f'(x)-\tfrac{1}{2}yyf'''(x)\ldots}
     = y\frac{f''(x)}{f'(x)}\Bigl\{1-yy\Bigl(\tfrac{1}{6}\frac{f^{\mathrm{IV}}(x)}{f''(x)}
                                  -\tfrac{1}{2}\frac{f'''(x)}{f'(x)}\Bigr)\ldots\Bigr\}.
\]
Therefore, by p.~144,
\[
   \tang c = D\cdot\frac{\sqrt{(1-ee\sin\varphi^{2})}}{a}\,y\,\tang\varphi,
\]
where
\begin{align*}
   \loghyp D
     &= \Bigl(-\tfrac{1}{6}\frac{f^{\mathrm{IV}}(x)}{f''(x)}
                + \tfrac{1}{2}\frac{f'''(x)}{f'(x)}\Bigr)yy\ldots\\
     &= -\tfrac{1}{2}yy\,\frac{(1-ee\sin\varphi^{2})^{3}}{aa(1-ee)^{3}}\,
         (1-3ee+2ee\sin\varphi^{2})\ldots\\
     &\text{or, in the approximation,}\\
     &= -\tfrac{1}{2}\frac{yy}{aa}(1-ee).
\end{align*}
Furthermore,
\[
   c = y\frac{f''(x)}{f'(x)}
      - \tfrac{1}{3}y^{3}\Bigl\{\frac{f^{\mathrm{IV}}(x)}{f'(x)}
                                - \frac{3f''(x)f'''(x)}{(f'(x))^{2}}
                                + 2\Bigl(\frac{f''(x)}{f'(x)}\Bigr)^{3}\Bigr\}\ldots
\]
or, abbreviating $f''(x)/f'(x) = h(x)$,
\[
   c = y\,h(x) - \tfrac{1}{3}y^{3}\,h''(x)\ldots.
\]

The associated rational functions in $\varphi$ are
\begin{align*}
   h(x) &= \frac{(1-ee\sin\varphi^{2})^{3/2}\sin\varphi}{a\cos\varphi},\\[2pt]
   h'(x) &= \frac{1-ee\sin\varphi^{2}}{aa(1-ee)\cos\varphi^{2}}
              \bigl\{1-2ee\sin\varphi^{2}+ee\sin\varphi^{4}\bigr\},\\[2pt]
   h''(x) &= \frac{2(1-ee\sin\varphi^{2})^{3/2}\sin\varphi}{a^{3}(1-ee)^{2}\cos\varphi^{4}}
                \bigl\{1-3ee+(2ee+4e^{4})\sin\varphi^{2}
                        -(ee+5e^{4})\sin\varphi^{4}+2e^{4}\sin\varphi^{6}\bigr\}\\[1pt]
            &= \frac{2(1-ee\sin\varphi^{2})^{3/2}\sin\varphi}{a^{3}(1-ee)^{2}\cos\varphi^{4}}
                \bigl\{(1-ee)^{2}-ee(1-ee)\cos\varphi^{4}-2e^{4}\cos\varphi^{6}\bigr\}.
\end{align*}

Therefore
\begin{multline*}
   c = \frac{\sqrt{(1-ee\sin\varphi^{2})}}{a}\,\tang\varphi\cdot y
   - \tfrac{1}{3}\,\frac{(1-ee\sin\varphi^{2})^{3/2}\sin\varphi}{a^{3}(1-ee)^{2}\cos\varphi^{3}}
       \bigl\{1-3ee+(2ee+4e^{4})\sin\varphi^{2}\\
   -(ee+5e^{4})\sin\varphi^{4}+2e^{4}\sin\varphi^{6}\bigr\}y^{3}\ldots
\end{multline*}
or, in the alternative form,
\begin{multline*}
   c = \frac{\sqrt{(1-ee\sin\varphi^{2})}}{a}\,\tang\varphi\cdot y\\
   - \tfrac{1}{3}\,\frac{(1-ee\sin\varphi^{2})^{3/2}\sin\varphi}{a^{3}(1-ee)^{2}\cos\varphi^{3}}
       \bigl\{(1-ee)^{2}-ee(1-ee)\cos\varphi^{4}-2e^{4}\cos\varphi^{6}\bigr\}y^{3}\ldots
\end{multline*}
Introducing $\lambda$ as the principal independent variable one
also finds
\[
   c = \lambda\sin\varphi
     - \frac{(1-ee\sin\varphi^{2})^{3/2}\tang\varphi}{6a^{3}(1-ee)^{3}}
        \bigl\{(1-ee)^{2}-3ee(1-ee)\cos\varphi^{2}-4e^{4}\cos\varphi^{4}\bigr\}y^{3}\ldots.
\]

\bigskip

\subsection*{[3.] (scan p.~148, glosses \texttt{[wo $H=\tfrac12 k\rho^{2}\cdot10^{7}$, $\log H = 5{,}531\,8128-10$]} expanded)}

For numerical computation the most convenient formulae are obtained
by setting
\begin{align*}
   L &= \frac{(1-ee\sin\varphi^{2})^{3/2}\,y}{a\rho\cos\varphi},\\[2pt]
   M &= \frac{(1-ee\sin\varphi^{2})^{2}\,\tang\varphi\cdot yy}{2aa(1-ee)\rho},\\[2pt]
   N &= \frac{(1-ee\sin\varphi^{2})^{3/2}}{a\rho}\,\tang\varphi\cdot y.
\end{align*}
Then
\begin{align*}
   \log\lambda &= \log L - A\cdot LL,\\
   \log(\varphi-\Phi) &= \log M - B\cdot LL,\\
   \log c &= \log N - C\cdot LL,
\end{align*}
where the logarithms of $A,B,C$ are most conveniently taken from a
separate table with argument $\varphi$.
For computing this table use
\begin{align*}
   A &= H\Bigl\{0{,}75 - \tfrac{1}{4}\cos 2\varphi
           + \tfrac{ee}{2(1-ee)}\cos\varphi^{4}\Bigr\},\\
   B &= H\Bigl\{0{,}75
           + \tfrac{2-11ee}{4(1-ee)}\cos\varphi^{2}
           + \tfrac{5ee}{2(1-ee)}\cos\varphi^{4}
           - \tfrac{e^{4}}{(1-ee)^{2}}\cos\varphi^{6}\Bigr\},\\
   C &= H\Bigl\{1 - \tfrac{ee}{1-ee}\cos\varphi^{4}
           - \tfrac{2e^{4}}{(1-ee)^{2}}\cos\varphi^{6}\Bigr\},
\end{align*}
where (the gloss explicitly):
\[
   H = \tfrac{1}{2}k\rho\rho\cdot 10^{7},\qquad
   \log H = 5{,}531\,8128-10,
\]
with $k$ the modulus of Briggsian (common) logarithms and
$\rho = 1/206264{,}8\ldots$.

The corrections $A\cdot LL$, $B\cdot LL$, $C\cdot LL$ are obtained in
units of the seventh decimal; $\lambda$, $\varphi-\Phi$ and $c$ come
out in seconds. The subsequent table of $\log A$, $\log B$, $\log C$
in the source TeX is based on the flattening $1/302{,}68$ (the
Hannover value); cf.\ scan p.~149.

\bigskip
\hrule
\bigskip

\section*{Audit cross-references}

\begin{itemize}\setlength{\itemsep}{2pt}
\item \texttt{gauss\_b01a\_ch4.tex} lines 480--530 ([6.]--[9.]):
      expanded above (Sects.\ [6.]--[9.] of this file).
\item Source TeX line 512--517 (the explicit 5-line stub of [9.]):
      now Section [9.] of this file. The whole Laplace-type identity
      plus its integration to a series in $y$ is restored.
\item Source TeX line 397 ([Tabular data for $\log A$, $\log B$, $\log C$
      follows --- transcribed as placeholder]): the table itself is in
      the Hannover scan p.~149 and is referenced but not retyped here
      (the audit explicitly noted this as a known gap for Phase~2
      manual table re-entry).
\item Source TeX line 487--491 ([6.]--[7.] ``Weitere Entwicklungen f\"ur
      die Coordinatenberechnung''): expanded above to scanned content.
\item Source TeX line 493--510 ([8.]--[9.] ``Allgemeine Formeln f\"ur
      die Differentiale''): fully written out as Sect.~[8.]
      (definitions of $\alpha,\beta,\gamma,\delta,\varepsilon$) and
      Sect.~[9.] of this file (the Laplace-type identity).
\item Source TeX line 529--536 ([10.]--[14.] ``Entwicklung der
      Grundformeln \dots'' --- summarised in source as [Die Entwicklung
      beruht auf folgender Aufgabe.]): the auxiliary lemma is in the
      source TeX correctly; the surrounding [Schlussbemerkungen und
      Tabellen] stub is intentionally left as-is (numerical tables
      only).
\end{itemize}

\bigskip
\hrule
\bigskip

\noindent\hfill\textsc{Expanded against scan, end of file.}

\end{document}
